% Article 1 — Leakage-audited benchmark on GeNIS 2025
% Target: Journal of Information Security and Applications (elsarticle)
% Build: pdflatex main && bibtex main && pdflatex main && pdflatex main
\documentclass[review,12pt]{elsarticle}

\usepackage{amsmath,amssymb}
\usepackage{booktabs}
\usepackage{pifont}
\usepackage{graphicx}
\usepackage{hyperref}
\usepackage{xcolor}

% Working notes to co-authors: \todo{...} prints in red; remove before submission.
\newcommand{\todo}[1]{\textcolor{red}{[TODO: #1]}}

\journal{Journal of Information Security and Applications}

\begin{document}

\begin{frontmatter}

\title{A Leakage-Audited Benchmark of Deep and Ensemble Detectors on the
GeNIS 2025 Corpus: Calibration, Cost, and Class Imbalance under the
Natural Distribution}

\author[1]{Ala Bahri}
\author[1]{Farah Jemili}
\author[2]{Mohamed Mosbah}
\affiliation[1]{organization={\todo{affiliation 1}}}
\affiliation[2]{organization={\todo{affiliation 2}}}

\begin{abstract}
\todo{Write last. Structure to follow: (1) modern corpora are needed but must be
audited before they become reference numbers; (2) we benchmark 9 detector
families on GeNIS 2025 under a clean, natural-distribution protocol;
(3) headline finding: the shortcut audit --- a timestamp-only classifier
reaches 99.3\% (chance 11\%), and the model ranking changes materially once
positional and quasi-identifying features are excluded and a chronological
split is imposed; (4) interval study: how 5/10/30/60\,s aggregation changes
per-family detectability and cost; (5) deliverables: per-interval feature
blacklist, calibrated operating points, CPU inference cost, full artefact.}
\end{abstract}

\begin{keyword}
intrusion detection \sep GeNIS \sep benchmark \sep data leakage \sep
shortcut learning \sep model calibration \sep class imbalance \sep
inference cost
\end{keyword}

\end{frontmatter}

% =====================================================================
\section{Introduction}
\label{sec:intro}

Learning-based network intrusion detection is evaluated, almost universally,
on a small set of shared corpora, and the field has learned the hard way that
the corpus shapes the conclusions. CICIDS2017, the de facto reference for
nearly a decade, was later shown to carry feature-extraction and labelling
errors that materially change reported performance \todo{cite Engelen 2021,
Liu 2022, Lanvin 2022}; systems tuned to its idiosyncrasies aged with it. As
the community moves to a new generation of corpora, the question is not only
\emph{which detectors perform best}, but \emph{whether the first published
numbers on a new corpus measure traffic patterns or artefacts of its
generation}. History suggests the audit should come before the leaderboard,
not a decade after.

GeNIS \todo{cite Silva DIB 2025} is one of the most promising corpora of this
new generation: captured in February 2025 by an independent laboratory (GECAD,
Polytechnic of Porto) on an Airbus CyberRange simulation of a small and
medium-sized enterprise network, exported with the HERA flow exporter at four
aggregation intervals (5, 10, 30 and 60 seconds, 2.8 million flows in total),
and publicly archived. The companion study by its authors \todo{cite Silva
FPS 2025} validates binary and multiclass classification with five feature
selection methods over tree ensembles and deep networks, and constitutes the
official baseline. What no published work provides, to our knowledge, is an
\emph{independent} evaluation that asks the auditing questions: which features
carry shortcuts, what happens under a chronological split, how calibrated the
resulting confidences are, what the detectors cost at inference time, and how
the choice of aggregation interval changes the task.

This paper provides that evaluation. Its starting observation, which motivates
everything downstream, is that on the 60-second slice of GeNIS \emph{a
classifier given only the capture timestamp reaches 99.3\% nine-class
accuracy, against a chance level of 11\%}: attack windows are narrow and
mutually disjoint in capture time, so any feature correlated with position in
the capture is a shortcut, and any random-split evaluation is at risk of
rewarding memorisation of the capture schedule rather than recognition of
traffic behaviour. \todo{Verify the 99.3\% figure against the E3 re-run on
this paper's frozen protocol before submission.}

We therefore benchmark nine model families --- three deep detectors (RNN, 1D
CNN, DNN), three boosted-tree and forest ensembles (Random Forest, XGBoost,
LightGBM), a modern tabular transformer (FT-Transformer), an
unsupervised autoencoder evaluated separately as an unknown-attack detector,
and trivial baselines --- under a frozen, leakage-audited protocol: scaler
fitted on the training split only, natural class distribution preserved (7.4\%
benign traffic in the 60-second test partition), five seeds, and paired
statistical tests. On top of this common protocol we run the audit
(permutation importance, stratified-versus-chronological comparison, feature
exclusion), measure calibration before and after temperature scaling, meter
CPU inference cost, and repeat the study across the four aggregation
intervals.

\paragraph{Contributions}
\begin{enumerate}
  \item \textbf{A shortcut and leakage audit of GeNIS}, quantifying how much
  of the near-perfect random-split performance is attributable to positional
  and quasi-identifying features, and delivering a per-interval feature
  blacklist as a reusable artefact for any future evaluation on this corpus.
  \item \textbf{The first calibrated, cost-aware, imbalance-aware benchmark}
  of deep and ensemble detectors on GeNIS under the natural distribution,
  with macro-F1, MCC, PR-AUC, operational FPR, ECE and reliability
  diagrams, and CPU throughput/latency for every model.
  \item \textbf{The first cross-interval study} of the corpus, measuring how
  5/10/30/60-second flow aggregation shifts per-family detectability and
  inference cost.
  \item \textbf{A public, single-command-per-experiment artefact} (code,
  configurations, seeds, blacklist, checkpoints) intended as the reference
  protocol for subsequent work on GeNIS.
\end{enumerate}

\todo{Closing paragraph: paper organisation, once section numbering is stable.}

% =====================================================================
\section{Related work}
\label{sec:related}

\subsection{Benchmark corpora and their documented flaws}
Doubts about how faithfully laboratory corpora measure intrusion detection
predate deep learning: Sommer and Paxson's classic analysis already warned
that closed-world evaluations reward properties of the data rather than of
the detector \cite{sommer2010closed}. The subsequent decade validated the
warning on the field's main reference. CICIDS2017 \cite{sharafaldin2018cicids} became the de facto benchmark, and was then shown to carry defects at
every level of its production: errors in the CICFlowMeter feature extractor
and in the labels themselves \cite{engelen2021troubleshooting}, an error
prevalence material enough to change reported performance
\cite{liu2022error}, and further labelling mistakes whose correction
significantly alters detection results \cite{lanvin2022errors}. Meanwhile, the
most-cited benchmark on that corpus \cite{maseer2021benchmarking} ---
ten supervised and unsupervised algorithms compared on accuracy, F1 and
training/testing time --- predates the shortcut-learning literature and
evaluates under random splits without a leakage or temporal audit, so its
reference numbers inherit whatever artefacts the corpus carries. The lesson
is not that testbed corpora or model-zoo benchmarks are useless; it is that
a corpus's first years of use fix the community's reference numbers, so the
audit of a new corpus is most valuable \emph{before} a literature
accumulates on it. GeNIS
\cite{silva2025genis} is exactly at that stage: captured in February 2025 on
an Airbus CyberRange simulation of a small and medium-sized enterprise
network, exported at four aggregation intervals with the HERA tool
\cite{pinto2024hera}, built on the Argus flow exporter, and released with
raw packets so that alternative feature extractions remain possible --- a
design choice motivated by the observation that the exporter itself
materially affects downstream model performance \cite{pinto2025exporter}.
Its authors' companion study \cite{silva2025binary} validates binary and
multiclass classification with five feature-selection methods over tree
ensembles and deep networks, and stands as the official baseline this paper
is anchored to.

\subsection{Shortcut learning and evaluation bias in learned security}
That near-perfect benchmark scores can rest on features unrelated to the
phenomenon of interest is now recognised across machine learning as
\emph{shortcut learning} \cite{geirhos2020shortcut}, and security is a
domain where the stakes of the confusion are highest. Arp et al.\ catalogue
the recurring pitfalls of security ML --- sampling bias, spurious
correlations, data snooping --- and find them endemic in published work
\cite{arp2022dodo}. Two lines are directly upstream of ours. TESSERACT
demonstrated that ignoring temporal structure inflates malware-classification
results, and argued for chronologically consistent splits as a validity
requirement rather than a robustness extra \cite{pendlebury2019tesseract}.
Closer to network flows, D'hooge et al.\ established the contaminating
effect of metadata features in NIDS models: the destination port alone
separates ten state-of-the-art IDS datasets at 70--100\% accuracy, so any
model given that feature learns a dataset artefact rather than traffic
behaviour \cite{dhooge2022metadata}. Our audit extends this programme along
three axes: from metadata to \emph{positional} features (capture timestamps
and exporter sequence counters, which on GeNIS predict the class at
$\approx$99\% against an 11\% chance level \todo{confirm E3 figure}); from a
single split convention to a stratified-versus-chronological comparison that
measures how much of the ranking survives; and from the CICIDS family to a
new-generation corpus whose reference numbers are not yet fixed. The
per-interval feature blacklist we release is the practical form of
D'hooge et al.'s countermeasure argument, instantiated for GeNIS.

\subsection{Calibration, imbalance, and the cost of a verdict}
A benchmark meant to inform deployment must report more than a ranking.
Modern neural networks are systematically over-confident
\cite{guo2017calibration}, and a detector whose confidence is dishonest
cannot drive thresholded automation; temperature scaling
\cite{guo2017calibration} and selective classification with a reject option
\cite{geifman2017selective,hendrickx2024reject} are the standard correctives,
and uncertainty quantification is increasingly argued to be a first-class
requirement for trustworthy NIDS \cite{talpini2024uncertainty}. Class
imbalance interacts with all of this: under the natural distribution of
GeNIS (7.4\% benign traffic at the 60-second interval), accuracy is close to
vacuous and metrics such as macro-F1 and the Matthews correlation
coefficient are the informative ones \cite{chicco2020mcc}. Finally, the
tabular-learning literature gives strong prior reason to include tree
ensembles alongside deep detectors: gradient-boosted trees remain
consistently competitive or superior on tabular data
\cite{grinsztajn2022trees,shwartzziv2022tabular}, with attention-based
architectures such as the FT-Transformer as the strongest deep contenders
\cite{gorishniy2021ftt}. Inference cost is one dimension the GeNIS baseline
already reports \cite{silva2025binary}; calibration, selective operating
points, and confidence-bearing comparison are not, and this paper adds them
to the same protocol.

\subsection{Prior work on GeNIS and positioning}
At the time of writing, the published record on GeNIS consists of the
dataset descriptor \cite{silva2025genis}, its authors' classification
baseline \cite{silva2025binary}, and one external study that uses the corpus
for machine unlearning \cite{magalhaes2026unlearning}, a goal orthogonal to
benchmarking.

The authors' baseline deserves a precise description, because this paper
builds on it rather than competing with it, and several of its choices are
already careful ones. Silva et al.\ combine five feature-selection methods
(information gain, chi-squared, recursive feature elimination, mean absolute
deviation, dispersion ratio) into a single ranking, retain the 16
highest-scoring features, and train three tree ensembles (RF, XGB, LGBM) and
two neural networks (LSTM, MLP) with grid-searched hyperparameters under
5-fold cross-validation. They evaluate under the natural class distribution
(7.37\% benign), report training, per-epoch and inference times alongside
detection metrics, deliberately exclude topology-bound features --- IP, MAC
and VLAN identifiers, and the HERA-specific \texttt{Ssaddr}/\texttt{Sdaddr}
counters --- and use SHAP to check which feature families drive predictions.
They close by flagging a concern they do not pursue: the models show an
``overreliance on traffic volume or packet counts''.

Four properties of that study define the space this paper occupies. First,
its multiclass task uses the coarse \texttt{CategoryLabel} taxonomy --- four
classes (benign, DoS, reconnaissance, brute force) --- whereas we evaluate
the nine-class \texttt{SubCategoryLabel} task, which requires separating five
denial-of-service families from one another rather than pooling them.
Second, evaluation rests on a single stratified split with no repetition, so
no confidence interval or significance test accompanies the reported
differences. Third, only the 60-second interval is used, leaving the effect
of flow aggregation unmeasured. Fourth, and most consequentially for a corpus
whose attack windows are disjoint in capture time
(Section~\ref{sec:chronology}), the protocol is a random split with no
temporal evaluation and no test for whether retained features encode capture
position rather than traffic behaviour --- the audit that this paper
supplies, and which turns their own SHAP-based concern into a measurement.
Table~\ref{tab:positioning} situates these dimensions.

\begin{table}[t]
\centering
\caption{Positioning against the closest related work. \ding{51} covered,
$\pm$ partial, \ding{55} absent. ``New corpus'': a post-2024 corpus.
``Repeated'': multiple seeds with dispersion or significance testing.
``Cost'': measured training and inference time.}
\label{tab:positioning}
\small
\setlength{\tabcolsep}{4pt}
\begin{tabular}{lcccccccc}
\toprule
 & New & Shortcut & Temporal & Calib- & Natural & Repea- & Cost & Multi-\\
 & corpus & audit & protocol & ration & distrib. & ted & & interval\\
\midrule
Maseer et al.\ \cite{maseer2021benchmarking} & \ding{55} & \ding{55} & \ding{55} & \ding{55} & $\pm$ & \ding{55} & \ding{51} & \ding{55}\\
Silva et al.\ \cite{silva2025binary} & \ding{51} & \ding{55} & \ding{55} & \ding{55} & \ding{51} & \ding{55} & \ding{51} & \ding{55}\\
D'hooge et al.\ \cite{dhooge2022metadata} & \ding{55} & \ding{51} & \ding{55} & \ding{55} & \ding{55} & $\pm$ & \ding{55} & \ding{55}\\
TESSERACT \cite{pendlebury2019tesseract} & \ding{55} & $\pm$ & \ding{51} & \ding{55} & \ding{51} & \ding{51} & \ding{55} & \ding{55}\\
Uncertainty-aware NIDS \cite{talpini2024uncertainty} & \ding{55} & \ding{55} & \ding{55} & \ding{51} & $\pm$ & $\pm$ & \ding{55} & \ding{55}\\
\textbf{This work} & \ding{51} & \ding{51} & \ding{51} & \ding{51} & \ding{51} & \ding{51} & \ding{51} & \ding{51}\\
\bottomrule
\end{tabular}
\end{table}

% =====================================================================
\section{The GeNIS corpus and its natural distribution}
\label{sec:corpus}

\subsection{Origin and generation}
GeNIS (GECAD Network Intrusion Scenarios) \cite{silva2025genis} was recorded
between 6 and 12 February 2025 on the Airbus CyberRange platform, configured
as a digital twin of a small and medium-sized enterprise network: six
segmented LANs (administration, users, internal servers, a DMZ hosting
public-facing web/FTP/mail services, remote VPN users, and the attacker
hosts), with port mirroring feeding an isolated capture machine running
\texttt{dumpcap}. Benign traffic was generated by the Benign User Profiler
\cite{shafi2024bup} under three profiles --- office-user activity, system
administration, and passive background traffic --- while the malicious
traffic realises eight sequential attack scenarios, each following the same
kill-chain skeleton (DNS resolution, Nmap scanning, then the terminal
action): five denial-of-service floods against the DMZ web server (Hulk,
Slowloris via Metasploit, and UDP/ICMP/PSH-ACK floods via
\texttt{hping3}) and three credential brute-force attacks (SMB, SSH, and
combined SSH+FTP) using Hydra with a 10,000-password dictionary. The raw
packets (37.7 million) were aggregated into labelled flows by the HERA
exporter \cite{pinto2024hera}, built on Argus, at four flow intervals ---
5, 10, 30 and 60 seconds --- for a total of 2.8 million flows carrying
hierarchical labels (binary, category, subcategory).

\subsection{The slice used in this study}
The corpus is released in a modular layout; this benchmark uses the
\texttt{2-flows} module, which provides per-subcategory CSV files at all
four intervals. Two taxonomic points must be fixed before any comparison.
First, the reconnaissance subcategories (\texttt{recon-nmap},
\texttt{recon-dns}) listed in the corpus descriptor do not appear as
standalone files in \texttt{2-flows}; they exist only inside the
\texttt{3-scenarios} recompositions, so the classification task supported
by this module comprises \textbf{eight attack families}. Second, the three
benign profiles (user, administrator, background) are variants of normal
traffic rather than separate detection targets; we merge them into a single
\texttt{benign} class, yielding the nine-class task evaluated throughout
(benign, \texttt{bruteforce-\{ftp,smb,ssh\}},
\texttt{dos-\{hulk,icmp,pshack,slowloris,udp\}}).

These choices place our task strictly between the two published
alternatives, and the difference is worth stating precisely because it
governs comparability. The corpus baseline \cite{silva2025binary} evaluates
the four-class \texttt{CategoryLabel} task on 368{,}556 sixty-second flows,
a count that includes reconnaissance; the \texttt{2-flows} module we use
yields 338{,}820. The 29{,}736-flow gap is fully accounted for by two
independent causes, which we report because neither is documented: the
absence of the reconnaissance captures (\texttt{recon-nmap}, 27{,}713 flows;
\texttt{recon-dns}, 20) from this module, and a 2{,}003-flow shortfall in
\texttt{benign-background} relative to the descriptor's Table 12
(8{,}283 released versus 10{,}286 reported). All figures in this paper are
recounted from the released files, and we state the released counts rather
than the documented ones wherever the two disagree.
\todo{Insert the recomputed per-interval statistics table (Table 1 from the
notebook run) and confirm the two counts above against it.}

\subsection{Capture chronology and its consequences}
\label{sec:chronology}
The capture schedule matters as much as the volumes: benign user and
administrator activity was recorded on Thursday 6 and Friday 7 February,
background traffic over the weekend of 8 February, and \emph{all} attacks
during the following week, 10--12 February. Benign and malicious traffic
are therefore \textbf{completely segregated in capture time}, and each
attack family occupies narrow, mutually disjoint windows within the attack
days (Figure~\ref{fig:timeline}). Three consequences structure the rest of
this paper. (i) Any feature correlated with position in the capture ---
timestamps, but also exporter counters --- is a class-revealing shortcut,
which motivates the audit of Section~\ref{sec:results}. (ii) A
\emph{global} chronological split is degenerate: the final portion of the
capture contains only attack blocks never seen in training, so every model
scores near zero --- a measurement of the schedule, not of generalisation
(we report it as such). (iii) Temporal evaluation must therefore be
\emph{per class}: our temporal protocol orders each class's flows in time
and tests on its later flows, keeping every family represented while
guaranteeing that every test flow postdates every training flow of its
class (Section~\ref{sec:protocol}).

\begin{figure}[t]
\centering
% figures/fig1_timeline.pdf from the notebook run
\todo{include fig1\_timeline}
\caption{Activity windows per class over the 6--12 February 2025 capture.
Benign traffic (green) is confined to the first days; each attack family
(red) occupies narrow, disjoint windows in the attack days.}
\label{fig:timeline}
\end{figure}

% =====================================================================
\section{Evaluation protocol}
\label{sec:protocol}
\todo{4.1 Clean stratified protocol (scaler on train only, 60/20/20, 5 seeds,
natural distribution). 4.2 Chronological protocol. 4.3 The audit procedure:
timestamp-only probe, permutation importance, exclusion rule, blacklist
construction --- three feature categories: positional/timestamp,
quasi-identifiers, behavioural-but-collapsing. 4.4 Hyperparameter budget:
$\sim$30 random-search trials AND equal compute ceiling per model, both
stated. 4.5 Metrics and statistics: macro-F1, MCC, PR-AUC, operational
FPR/FNR, ECE (15 bins); McNemar + Holm, bootstrap 95\% CIs (10k).}

% =====================================================================
\section{Models}
\label{sec:models}
\todo{Table of the 9 families with exact architectures/search spaces. Deep
detectors reuse the GeNIS architectures of the companion system paper
(Table 3 there). Autoencoder gets its own subsection stating the separate
evaluation regime (AUROC on reconstruction error; optional
leave-one-family-out).}

% =====================================================================
\section{Experiments and results}
\label{sec:results}
\todo{6.1 E1 anchored comparison with Silva et al. (state deviations).
6.2 E2 main benchmark, 60s stratified --- expect saturation; present as
anchor, not headline. 6.3 E3 audit: ranking before/after (headline figure).
6.4 E4 intervals: per-family detectability vs 5/10/30/60s + cost.
6.5 E5 calibration. 6.6 E6 CPU cost bench. 6.7 (optional) E7 top-k ablation.}

% =====================================================================
\section{Discussion}
\label{sec:discussion}
\todo{What the audit implies for future GeNIS papers; recommendation of an
operating model for SME deployment (RQ4); relation to the companion
architecture paper (these checkpoints are its provenance).}

% =====================================================================
\section{Threats to validity}
\label{sec:threats}
\todo{Internal: residual leakage; temporal split as designed result.
External: CyberRange testbed, SME profile, volumetric floods, coarse FPR
granularity at low benign share. Construction: finite budget; AE evaluated
apart.}

% =====================================================================
\section{Conclusion}
\label{sec:conclusion}
\todo{Write last.}

\section*{Data and code availability}
\todo{Repository URL + Zenodo DOI of the artefact; GeNIS is public
(doi:10.5281/zenodo.14919237).}

\bibliographystyle{elsarticle-num}
\bibliography{references}

\end{document}
